\begingroup
\setlength{\abovedisplayskip}{7pt}
\setlength{\belowdisplayskip}{7pt}
\setlength{\abovedisplayshortskip}{3pt}
\setlength{\belowdisplayshortskip}{4pt}
\def\ixrad{\mathord{\vcenter{\hbox{\Large\(\surd\)}}}}
\def\ixangle#1#2#3{%
  \makebox[1.95em][r]{\(#1\)}%
  \makebox[1.45em][r]{\(#2\)}%
  \makebox[3.10em][r]{\(#3\)}}
\def\ixobsCount#1#2{\makebox[.8em][r]{#1}\enspace\makebox[2.6em][c]{#2}}
\def\ixobs#1#2#3#4#5#6#7#8#9{%
  #1 & \(\left\{\begin{array}{@{}l@{\enspace}l@{}}
  \text{Direkt}&\ixangle{#2}{#3}{#4}\\
  \text{Reflektiert}&\ixangle{#5}{#6}{#7}
  \end{array}\right.\) &
  \begin{tabular}{@{}l@{}}#8\\#9\end{tabular}}

% Printed 32; principal PDF 39; Google PDF 43.
\par\clearpage
\noindent\makebox[12mm][l]{\fontsize{13}{16}\selectfont 32}\parbox[t]{\dimexpr\linewidth-24mm\relax}{\centering\scriptsize\bfseries BESTIMMUNG DES BREITENUNTERSCHIEDS ZWISCHEN GÖTTINGEN UND ALTONA.}\makebox[12mm][r]{}\par\medskip

{\small
\setlength{\arraycolsep}{.85em}
\[
\begin{array}{r||r|r||r|r}
\smash{\raisebox{-.7ex}{Stern}}&
\multicolumn{2}{c||}{\text{\normalsize\textls[100]{Göttingen.}}}&
\multicolumn{2}{c}{\text{\normalsize\textls[100]{Altona.}}}\\
&\text{Koll. F.}&\text{Gewicht}&\text{Koll. F.}&\text{Gewicht}\\
\hline
24&3''{,}93&10{,}91&1''{,}21&12{,}00\\
25&3{,}88&12{,}00&0{,}83&12{,}00\\
26&3{,}54&12{,}00&0{,}76&12{,}00\\
27&4{,}44&4{,}00&1{,}97&7{,}50\\
28&3{,}27&10{,}91&0{,}95&12{,}00\\
29&3{,}89&6{,}00&1{,}15&12{,}00\\
30&4{,}54&10{,}91&1{,}70&12{,}00\\
31&4{,}52&10{,}91&1{,}73&12{,}00\\
32&3{,}26&12{,}00&1{,}13&12{,}00\\
33&3{,}49&10{,}91&1{,}11&10{,}91\\
34&4{,}58&12{,}00&1{,}66&12{,}00\\
35&3{,}61&10{,}00&1{,}37&10{,}00\\
36&4{,}33&12{,}00&1{,}34&10{,}00\\
37&3{,}97&12{,}00&1{,}91&10{,}91\\
38&3{,}52&12{,}00&1{,}56&10{,}00\\
39&4{,}16&3{,}00&1{,}45&8{,}00\\
40&5{,}04&4{,}80&1{,}02&6{,}00\\
41&3{,}46&4{,}80&1{,}75&10{,}00\\
42&3{,}84&4{,}00&2{,}45&8{,}89\\
43&4{,}40&4{,}00&1{,}01&10{,}00
\end{array}
\]
\par
}

Die Mittelwerte sind folgende:
\begin{center}
\begin{tabular}{@{}l@{\qquad}r@{\enspace}l@{}}
Kollimationsfehler in Göttingen & \(3''{,}75\) & mit dem Gewicht \(455{,}17\)\\
Kollimationsfehler in Altona & \(1{,}40\) & mit dem Gewicht \(432{,}18\).
\end{tabular}
\end{center}

Dass sich der Kollimationsfehler tatsächlich verändert hat, ist offensichtlich,
und es besteht kein Zweifel, dass dies während des Transports geschehen ist,
obwohl er mit aller möglichen Vorsicht durchgeführt wurde.

\begin{center}\textbf{4.}\end{center}

Obwohl man sich mit dem für den Breitenunterschied gefundenen
Ergebnis völlig zufriedengeben kann, ist es doch zumindest in theoretischer
Hinsicht sinnvoll, darauf hinzuweisen, dass die im 1. Art. angewandte
Kombination der Beobachtungen noch nicht die vorteilhafteste ist,
sofern nicht an jedem Ort jeder Stern in der einen Lage des Sektors ebenso
oft beobachtet wurde wie in der anderen. Tatsächlich hat die Bestimmung der
wahren Zenitdistanz in Göttingen durch die Formel \(\tfrac12(a+a')\) das Gewicht
\(\dfrac{4\alpha\alpha'}{\alpha+\alpha'}\). Wäre nun der Kollimationsfehler in Göttingen genau
bekannt und \(=f\),
so würde die Bestimmung der wahren Zenitdistanz dort durch die Formel
\[
\frac{\alpha(a+f)+\alpha'(a'-f)}{\alpha+\alpha'}
\]

% Printed 33; principal PDF 40; Google PDF 44.
\par\clearpage
\noindent\makebox[12mm][l]{}\makebox[\dimexpr\linewidth-24mm\relax][c]{\small\textbf{III.\quad ERGEBNISSE.}}\makebox[12mm][r]{\fontsize{13}{16}\selectfont 33}\par\medskip

\noindent das Gewicht \(\alpha+\alpha'=\dfrac{4\alpha\alpha'}{\alpha+\alpha'}+
\dfrac{(\alpha-\alpha')^2}{\alpha+\alpha'}\) haben, d.\ h. ein größeres als nach der
anderen Methode, sooft \(\alpha\) und \(\alpha'\) ungleich sind. Ebenso verhält es sich mit
der wahren Zenitdistanz in Altona, und auf diese Weise würden selbst
Beobachtungen in nur einer Lage (wie die von Nr. 5) zur Steigerung der
Genauigkeit beitragen, wenn auch nur geringfügig. Die Kollimationsfehler
an beiden Orten sind zwar nicht mit \textls[100]{absoluter} Genauigkeit bekannt,
doch man überzeugt sich leicht, dass die Verwendung der für sie gefundenen
Mittelwerte das Gewicht nur ganz unbedeutend vermindert.

\begin{center}\textbf{5.}\end{center}

Will man jedoch ein reines Ergebnis erhalten, das den Anforderungen der
strengen Theorie vollständig genügt, so muss man die Bestimmung des
Breitenunterschieds, der Kollimationsfehler und der wahren Zenitdistanzen der
einzelnen Sterne an einem der beiden Orte als ein einziges Problem behandeln. Dabei
müssen diese unbekannten Größen (in unserem Fall 46 an der Zahl) aus allen
durch sie bestimmten beobachteten Größen (171) mithilfe ebenso vieler
Gleichungen abgeleitet werden, indem diese nach den Regeln der
Wahrscheinlichkeitsrechnung kombiniert werden. Setzt man die Kollimationsfehler
in Göttingen und Altona \(=f\) und \(g\), den Breitenunterschied \(=h\), die
wahre Zenitdistanz eines Sterns in Göttingen \(=k\), so ergeben sich aus den
Beobachtungen dieses Sterns die vier Gleichungen mit den Gewichten \(\alpha,\alpha',\beta,\beta'\):
\[
\begin{aligned}
a&=k-f\\
a'&=k+f\\
b&=k-g-h\\
b'&=k+g-h.
\end{aligned}
\]

Es ist kaum nötig, daran zu erinnern, dass es zur Erleichterung der Rechnung
vorteilhafter ist, statt jener unbekannten Größen die noch erforderlichen
Korrekturen einzuführen, die an den bereits recht genau bestimmten Werten
anzubringen sind. Bezeichnen die Zeichen \(f^0,g^0,h^0,k^0\) diese Näherungswerte,
so kann man annehmen
\[
k^0=\frac{\alpha(a+f^0)+\alpha'(a'-f^0)+\beta(b+g^0+h^0)+\beta'(b'-g^0+h^0)}
{\alpha+\alpha'+\beta+\beta'}.
\]

Bei Befolgung dieser Regel (die man bei Anwendung der Methode
\par\smallskip
\noindent\textbf{IX.}\hfill 5\hspace{.12\textwidth}\par

% Printed 34; principal PDF 41; Google PDF 45.
\par\clearpage
\noindent\makebox[12mm][l]{\fontsize{13}{16}\selectfont 34}\parbox[t]{\dimexpr\linewidth-24mm\relax}{\centering\scriptsize\bfseries BESTIMMUNG DES BREITENUNTERSCHIEDS ZWISCHEN GÖTTINGEN UND ALTONA.}\makebox[12mm][r]{}\par\medskip

\noindent der kleinsten Quadrate auf auch nur etwas komplizierte Fälle niemals
aus den Augen verlieren sollte) und Verwendung einer geeigneten \textls[100]{indirekten}
Lösungsmethode wird eine Arbeit, die ohne diese Regel und bei direkter
Elimination unerträglich langwierig wäre, zu einem leichten Spiel.

\begin{center}\textbf{6.}\end{center}

Aus dieser Rechnung, die hier nicht ausführlich dargestellt werden muss,
ergibt sich, dass die früheren Bestimmungen keine merkliche Korrektur erfahren.
Die Korrektur des Breitenunterschieds ist \(=-0''{,}014\),
die Korrektur des Kollimationsfehlers in Göttingen \(=+0''{,}012\), die
Korrektur des Kollimationsfehlers in Altona \(=-0''{,}014\). Daraus folgen die neuen
Bestimmungen
\begin{center}
\begin{tabular}{@{}l@{\quad}r@{\enspace}c@{\enspace}l@{}}
Breitenunterschied & \(2^\circ0'56''{,}51\), & Gewicht & \(=217{,}67\)\\
Kollimationsfehler in Göttingen & \(3{,}76\) & » & \(\phantom{=}457{,}03\)\\
Kollimationsfehler in Altona & \(1{,}39\) & » & \(\phantom{=}437{,}64\).
\end{tabular}
\end{center}

Die Änderungen der nach der Regel des vorhergehenden Artikels
zugrunde gelegten wahren Zenitdistanzen der einzelnen Sterne in Göttingen
liegen ebenfalls fast alle unter \(0''{,}01\). Die sich ergebenden Werte hier
aufzuführen, wäre überflüssig, da aus diesen Werten die oben
mitgeteilten Deklinationen der Sterne unter der Annahme abgeleitet wurden,
dass die Polhöhe des Beobachtungsorts \(51^\circ31'47''{,}92\) beträgt. Dagegen geben wir hier die
Unterschiede an, die nach Einsetzen der gefundenen Werte in die
171 Gleichungen übrig bleiben.

{\small
\[
\begin{array}{@{}c|c||c|c||c|c||c|c@{}}
\text{Stern}&\text{Unterschied}&\text{Stern}&\text{Unterschied}&
\text{Stern}&\text{Unterschied}&\text{Stern}&\text{Unterschied}\\[.8ex]
1&+0''{,}07&5&+0''{,}03&9&-0''{,}06&13&+0''{,}87\\
&+0{,}09&&-0{,}03&&-0{,}09&&+0{,}30\\
&-0{,}26&&\text{—}&&-0{,}23&&-1{,}35\\
&+0{,}13&&-0{,}03&&+0{,}97&&-0{,}04\\
2&+0{,}56&6&+0{,}62&10&-0{,}71&14&+0{,}40\\
&-0{,}08&&+0{,}03&&-0{,}25&&+0{,}21\\
&-0{,}12&&-1{,}95&&+0{,}70&&-0{,}29\\
&-0{,}53&&-0{,}02&&+0{,}44&&-0{,}48\\
3&+0{,}48&7&-0{,}52&11&+0{,}12&15&-0{,}04\\
&+0{,}34&&+0{,}16&&-1{,}03&&+0{,}19\\
&-0{,}70&&-0{,}08&&+0{,}72&&-0{,}04\\
&-0{,}18&&+0{,}52&&+0{,}71&&-0{,}11\\
4&-0{,}35&8&-0{,}56&12&+0{,}52&16&-0{,}07\\
&-0{,}35&&+0{,}12&&-0{,}87&&-0{,}60\\
&+0{,}79&&+0{,}76&&+0{,}80&&-0{,}05\\
&-0{,}17&&-0{,}23&&-0{,}24&&+0{,}72
\end{array}
\]
\par
}

% Printed 35; principal PDF 42; Google PDF 46.
\par\clearpage
\noindent\makebox[12mm][l]{}\makebox[\dimexpr\linewidth-24mm\relax][c]{\small\textbf{III.\quad ERGEBNISSE.}}\makebox[12mm][r]{\fontsize{13}{16}\selectfont 35}\par\medskip

{\small
\[
\begin{array}{@{}c|c||c|c||c|c||c|c@{}}
\text{Stern}&\text{Unterschied}&\text{Stern}&\text{Unterschied}&
\text{Stern}&\text{Unterschied}&\text{Stern}&\text{Unterschied}\\[.8ex]
17&-0''{,}06&24&-0''{,}17&31&-0''{,}41&38&+0''{,}30\\
&-0{,}52&&+0{,}18&&+1{,}11&&-0{,}19\\
&+0{,}96&&+0{,}16&&-0{,}59&&-0{,}24\\
&-0{,}35&&-0{,}20&&+0{,}09&&+0{,}11\\
18&-0{,}23&25&-0{,}03&32&+0{,}14&39&+0{,}90\\
&-0{,}49&&+0{,}22&&-0{,}87&&+1{,}71\\
&+1{,}09&&+0{,}46&&+0{,}63&&-0{,}82\\
&-0{,}34&&-0{,}67&&+0{,}11&&-0{,}69\\
19&+0{,}32&26&+0{,}90&33&+0{,}19&40&-1{,}81\\
&-0{,}36&&+0{,}46&&-0{,}35&&+0{,}75\\
&+0{,}10&&-0{,}06&&+0{,}41&&+0{,}60\\
&+0{,}08&&-1{,}32&&-0{,}16&&-0{,}14\\
20&-0{,}06&27&-0{,}14&34&-0{,}48&41&+0{,}39\\
&-0{,}26&&+1{,}22&&+1{,}17&&-0{,}21\\
&+0{,}66&&-0{,}71&&-0{,}62&&-0{,}38\\
&-0{,}44&&+0{,}46&&-0{,}07&&+0{,}35\\
21&-0{,}22&28&+0{,}46&35&-0{,}08&42&+0{,}09\\
&-0{,}85&&-0{,}53&&-0{,}37&&+0{,}25\\
&+0{,}63&&+0{,}52&&+0{,}25&&-1{,}02\\
&+0{,}37&&-0{,}36&&+0{,}21&&+1{,}11\\
22&+0{,}77&29&-0{,}80&36&-0{,}14&43&-0{,}42\\
&-0{,}55&&-0{,}53&&+1{,}02&&+0{,}86\\
&-1{,}55&&+0{,}58&&-0{,}47&&+0{,}29\\
&+1{,}60&&+0{,}10&&-0{,}58&&-0{,}47\\
23&-0{,}80&30&-0{,}83&37&+0{,}11&&\\
&+0{,}67&&+0{,}74&&+0{,}54&&\\
&-0{,}03&&-0{,}21&&-0{,}93&&\\
&+0{,}03&&+0{,}42&&+0{,}11&&
\end{array}
\]
\par
}

\begin{center}\textbf{7.}\end{center}

Die Summe der Produkte aus den Quadraten dieser 171 Unterschiede und
der jeweils entsprechenden Anzahl der Beobachtungen ist \(=292{,}8249\). Nach
dem bereits angeführten Lehrsatz (Theoria Comb. Observ. Art. 38) gilt als
Näherungswert des mittleren Fehlers einer einzelnen Beobachtung die
Quadratwurzel aus dem Bruch, dessen Zähler jene Summe und dessen Nenner
der Überschuss der Anzahl der verglichenen Beobachtungsdaten über die Anzahl
der daraus nach der Methode der kleinsten Quadrate abgeleiteten unbekannten
Größen ist, in unserem Fall \(171-46=125\). Daraus ergibt sich der mittlere Fehler
\(=1''{,}5308\), der sich von dem im 2. Art. gefundenen nur wenig unterscheidet.
Der im Endergebnis für den Breitenunterschied zu erwartende mittlere Fehler
wäre demnach
\[
=\frac{1''{,}5308}{\surd 217{,}67}=0''{,}1038.
\]
\par\vfill
\noindent\hfill \(5^{*}\)\hspace{.12\textwidth}\par

% Printed 36; principal PDF 43; Google PDF 47.
\par\clearpage
\noindent\makebox[12mm][l]{\fontsize{13}{16}\selectfont 36}\parbox[t]{\dimexpr\linewidth-24mm\relax}{\centering\scriptsize\bfseries BESTIMMUNG DES BREITENUNTERSCHIEDS ZWISCHEN GÖTTINGEN UND ALTONA.}\makebox[12mm][r]{}\par\medskip

\begin{center}\textbf{8.}\end{center}

Bei den bisherigen Rechnungen wurde vorausgesetzt, dass sich alle Fehler
der verschiedenen Beobachtungen als völlig unabhängig voneinander oder
als rein zufällig betrachten lassen. Diese Voraussetzung ist aber offenbar
nicht ganz richtig, da sich alle \(\alpha\) Beobachtungen, die zur Bestimmung
eines \(a\) beigetragen haben, aufgrund der Bauart des Instruments auf ein und
denselben Teilungspunkt beziehen. Sie enthalten daher außer den eigentlichen
rein zufälligen Beobachtungsfehlern auch den Fehler der Teilung an diesem
Punkt. Dasselbe gilt für \(a'\), \(b\) und \(b'\). Die Teilungsfehler sind ihrerseits
unbekannte Größen, die in Bezug auf die einzelnen 171 Beobachtungsergebnisse
auch als rein zufällig und voneinander unabhängig betrachtet werden können,
da sich die wenigen Fälle vernachlässigen lassen, in denen sich mehrere
dieser Ergebnisse auf denselben Teilungspunkt bezogen. Dieser Umstand
erfordert eine Änderung der obigen Rechnungen,
obwohl sich die Ergebnisse für praktische
Zwecke letztlich gar nicht ändern.

Bezeichnet man den eigentlichen mittleren Beobachtungsfehler, der nur
auf Zufallsursachen unter Ausschluss der Teilungsfehler zurückgeht, mit \(m\),
und den mittleren Teilungsfehler mit \(\mu\), so ist der gesamte mittlere
Beobachtungsfehler \(=\surd(mm+\mu\mu)\) zu setzen und der mittlere Fehler eines
Mittelwerts aus \(\alpha\) Beobachtungen, die sich auf denselben Teilungspunkt beziehen,
\[
=\ixrad\left(\frac{mm}{\alpha}+\mu\mu\right),
\]
\noindent oder wenn wir \(\mu\mu=mm\theta\) setzen,
\[
=m\ixrad\left(\frac1\alpha+\theta\right).
\]

Wenn wir also das Gewicht einer Beobachtung \textls[100]{ohne Teilungsfehler}
als Einheit annehmen, wird das Gewicht von \(a\) nun
\[
=\frac{\alpha}{1+\alpha\theta}
\]
\noindent sein, und entsprechend die Gewichte von \(a'\), \(b\), \(b'\) in dieser Reihenfolge
\[
=\frac{\alpha'}{1+\alpha'\theta},\qquad
\frac{\beta}{1+\beta\theta},\qquad
\frac{\beta'}{1+\beta'\theta}.
\]
Bei der ersten Kombinationsmethode wird man daher das Gewicht des
\par

% Printed 37; principal PDF 44; Google PDF 48.
\par\clearpage
\noindent\makebox[12mm][l]{}\makebox[\dimexpr\linewidth-24mm\relax][c]{\small\textbf{III.\quad ERGEBNISSE.}}\makebox[12mm][r]{\fontsize{13}{16}\selectfont 37}\par\medskip

\noindent Ergebnisses für den Breitenunterschied aus den Beobachtungen eines Sterns,
wenn man den vorigen Ausdruck
\[
\frac4{\dfrac1\alpha+\dfrac1{\alpha'}+\dfrac1\beta+\dfrac1{\beta'}}=p
\]
\noindent setzt, jetzt
\[
=\frac p{1+p\theta}
\]
\noindent ansetzen und mit diesen Gewichten den Mittelwert aus den 42 Bestimmungen
bilden. Bei der zweiten Kombinationsmethode hingegen braucht man nur
jeder der 171 Gleichungen ein Gewicht zuzuordnen, das durch eine der
Formeln \(\dfrac{\alpha}{1+\alpha\theta}\) und\ so\ weiter bestimmt wird.

Offenbar können sich das Endergebnis selbst und der darin zu erwartende
mittlere Fehler nur dadurch ändern, dass die neuen Gewichte den früheren
nicht proportional sind. Bei der vorigen Methode wurden lediglich die
Ergebnisse der längeren Beobachtungsreihen etwas zu stark gewichtet. Die
Berücksichtigung der Teilungsfehler gleicht ihre Gewichte einander umso
mehr an, je größer die angenommenen Teilungsfehler sind. Bei Beobachtungen
mit einem Instrument, dessen Teilungsfehler
die eigentlichen Beobachtungsfehler bei Weitem überwögen,
könnte man daher nur alle Bestimmungen als gleich
zuverlässig betrachten.

\begin{center}\textbf{9.}\end{center}

Die beschriebenen Methoden bieten also keinerlei Schwierigkeiten, sobald
der Koeffizient \(\theta\) bekannt ist. Einen Näherungswert dafür, auf den
es hier begreiflicherweise nur ankommt, kann man auf indirektem
Wege ermitteln.

Wir bemerken zunächst, dass die Beobachtungen selbst ein Mittel bieten,
den eigentlichen mittleren Beobachtungsfehler \(m\) sehr zuverlässig zu
bestimmen. Tatsächlich zeigt er sich unabhängig vom Teilungsfehler in den
Unterschieden zwischen den \textls[100]{einzelnen} Werten, aus denen jedes
\(a\) (oder \(a'\), \(b\), \(b'\)) als Mittelwert gebildet ist, oder in ihren Abweichungen
von diesem Mittelwert. Wäre \(\alpha\) sehr groß, so ließe sich die Summe der Quadrate
dieser Abweichungen der einzelnen Werte von \(a\) vom Mittelwert als eine ge-
\par

% Printed 38; principal PDF 45; Google PDF 49.
\par\clearpage
\noindent\makebox[12mm][l]{\fontsize{13}{16}\selectfont 38}\parbox[t]{\dimexpr\linewidth-24mm\relax}{\centering\scriptsize\bfseries BESTIMMUNG DES BREITENUNTERSCHIEDS ZWISCHEN GÖTTINGEN UND ALTONA.}\makebox[12mm][r]{}\par\medskip

\noindent näherte Bestimmung von \((\alpha-1)mm\) ansehen. Eine solche einzelne
Bestimmung kann in unserem Fall, in dem \(\alpha\) nie größer als 7 ist, vom
richtigen Wert stark abweichen. Die Summe aller 171 Teilsummen
(für alle \(a\), \(a'\), \(b\), \(b'\) und für alle Sterne) darf aber nach den Grundsätzen
der Wahrscheinlichkeitsrechnung nur wenig abweichen von
\[
(\Sigma(\alpha-1)+\Sigma(\alpha'-1)+\Sigma(\beta-1)+\Sigma(\beta'-1))mm,
\]
\noindent in unserem Fall von \(728mm\). Wir haben für die Summe
der 171 Teilsummen
\[
=844{,}50
\]
\noindent gefunden, woraus sich für \(m\) der sehr zuverlässige Wert
\[
=1''{,}0770
\]
\noindent ergibt, der deutlich kleiner ist als der im 2. und 7. Artikel gefundene. Damit
bestätigt sich der Einfluss der Teilungsfehler vollständig. Wegen dieses Einflusses
konnten die früher berechneten Zahlen kein reines Ergebnis liefern.

\begin{center}\textbf{10.}\end{center}

Solange der mittlere Teilungsfehler nicht direkt bekannt ist, kann man
\(\theta\) nun indirekt so bestimmen, dass bei Anwendung
der ersten Methode nach dem Verfahren des Art. 2 oder bei Anwendung
der zweiten Methode nach dem Verfahren des Art. 7 der mittlere Fehler
einer Beobachtung, deren Gewicht als Einheit angenommen wurde, wieder
dem gefundenen Wert von \(m\) gleich ist.

Es schien sich jedoch nicht zu lohnen, solche Versuche so lange zu
wiederholen, bis eine vollkommene Übereinstimmung erreicht wäre.
Nachdem sich aus anderen Überlegungen ergeben hatte, dass der endgültige
Wert von \(\theta\) nur wenig von \(0{,}2\) abweichen könnte, schien es hinreichend,
diesen Wert allein der ersten Kombinationsmethode zugrunde zu legen.
Daraus ergab sich
\begin{center}
\begin{tabular}{@{}l@{\quad}l@{}}
Breitenunterschied & \(=2^\circ0'56''{,}50\)\\
Gewicht dieser Bestimmung & \(=104{,}29\)\\
\multicolumn{2}{@{}l@{}}{Mittlerer Fehler einer Beobachtung mit Einheitsgewicht,}
\end{tabular}
\end{center}
\[
=1''{,}131
\]

% Printed 39; principal PDF 46; Google PDF 50.
\par\clearpage
\noindent\makebox[12mm][l]{}\makebox[\dimexpr\linewidth-24mm\relax][c]{\small\textbf{III.\quad ERGEBNISSE.}}\makebox[12mm][r]{\fontsize{13}{16}\selectfont 39}\par\medskip

\noindent und daher der im obigen Endergebnis zu erwartende mittlere Fehler
\[
=0''{,}1108.
\]

Die Anwendung der zweiten Kombinationsmethode mit demselben Wert
von \(\theta\) hätte vermutlich eine noch bessere Übereinstimmung mit dem obigen
Wert von \(m\) ergeben, das Endergebnis für den Breitenunterschied
vielleicht um \(0''{,}01\) vermindert und das Gewicht dieser Bestimmung gewiss
ein wenig vergrößert. Es schien aber die Mühe nicht wert, die Rechnung
deshalb erneut durchzuführen. Man kann sich also an den gefundenen
Breitenunterschied \(2^\circ5'56''50\) halten und annehmen, dass sein Fehler wahrscheinlich
innerhalb der Grenzen \(\pm0''{,}07\) liegt.

\begin{center}\textbf{11.}\end{center}

Wenn wir den obigen Wert von \(\theta\) beibehalten, ergibt sich der mittlere
Teilungsfehler \(=m\surd\theta=0''{,}48\), sodass man den sogenannten wahrscheinlichen
Teilungsfehler der einzelnen Punkte \(=0''{,}32\) setzen kann. Offenbar
bezieht sich dies aber nur auf die unregelmäßigen Teilungsfehler, also auf
die Abweichungen der einzelnen Punkte von einer gedachten gleichförmigen
Teilung, die diesen Punkten so genau wie möglich angepasst ist. Deren
\textls[100]{absolute} Richtigkeit konnte hier eigentlich gar nicht in Betracht kommen.
Mit anderen Worten, das für den Breitenunterschied gefundene Ergebnis bezieht
sich mit der ihm zugeschriebenen Genauigkeit streng genommen nur auf mittlere
Sektorgrade und bleibt von deren absoluter Richtigkeit abhängig.
Das Instrument selbst bietet dem Astronomen kein unabhängiges Mittel,
\textls[100]{diese} zu prüfen. Bedenkt man jedoch, dass der Instrumentenmacher die Endpunkte
des Bogens mit äußerster Sorgfalt festgelegt hat und es hier nur um
einen kleinen Teil des gesamten Bogens geht, so wird man zugeben müssen,
dass diese Ursache die Unsicherheit des gefundenen Breitenunterschieds
nur sehr geringfügig vergrößern kann. Eine gewisse Kontrolle der absoluten
Richtigkeit der Teilung bieten auch die von mir am \textsc{Reichenbach}schen
Meridiankreis beobachteten Zenitdistanzen derselben 43 Sterne.
Ihre Unterschiede zu den am Sektor beobachteten Zenitdistanzen zeigen
bei einer Anordnung nach den Deklinationen
keine Spur von Regelmäßigkeit.

% Printed 40; principal PDF 47; Google PDF 51.
\par\clearpage
\noindent\makebox[12mm][l]{\fontsize{13}{16}\selectfont 40}\parbox[t]{\dimexpr\linewidth-24mm\relax}{\centering\scriptsize\bfseries BESTIMMUNG DES BREITENUNTERSCHIEDS ZWISCHEN GÖTTINGEN UND ALTONA.}\makebox[12mm][r]{}\par\medskip

\begin{center}\textbf{12.}\end{center}

Der Mittelpunkt des Sektors in Göttingen lag \(1{,}060\) Toisen
nördlich und \(7{,}595\) Toisen östlich vom Zentrum der Achse des \textsc{Reichenbach}schen
Meridiankreises. In Altona dagegen lag der Mittelpunkt des Sektors \(13{,}511\)
Toisen südlich und \(2{,}578\) westlich vom Mittelpunkt des dortigen
Meridiankreises. Die Reduktion des Breitenunterschieds zwischen den Standorten der
Sektoren auf den zwischen den Standorten der Meridiankreise beträgt daher für Göttingen \(0''{,}07\) und für Altona \(0''{,}85\).
Daraus ergibt sich der Breitenunterschied der Sternwarten von Göttingen und Altona
in Bezug auf die Standorte der \textsc{Reichenbach}schen Meridiankreise
\[
=2^\circ0'57''{,}42.
\]

\begin{center}\textbf{13.}\end{center}

Die absolute Polhöhe, die den oben angegebenen, aus den Zenitdistanzen
abgeleiteten Deklinationen der Sterne zugrunde liegt, beruht auf 89
Beobachtungen des Polarsterns am \textsc{Reichenbach}schen Meridiankreis, in beiden
Kulminationen, direkt und von einer Wasseroberfläche reflektiert. Da die
Beobachtungen aus dem Jahr 1824, die den größten Teil ausmachen, bisher noch nicht
veröffentlicht wurden, stelle ich hier sämtliche Beobachtungen zusammen
und bemerke nur, dass die direkte Einstellung meistens beim Erreichen des
zweiten, vierten (mittleren) und sechsten Fadens vorgenommen wurde, die Einstellung
des reflektierten Bildes dagegen beim Erreichen des ersten, dritten, fünften und
siebten Fadens. Angegeben ist hier der Mittelwert dieser auf den Zeitpunkt der
Kulmination reduzierten Zenitdistanzen, der nach \textsc{Bessel}s Tafeln nur um
die Refraktion korrigiert wurde, also noch den Kollimationsfehler und die Wirkung
der Biegung des Fernrohrs enthält.

\begin{center}
\textls[100]{Zenitdistanzen des Polarsterns.}\par\smallskip
1820.\quad Kreis im Osten.\par\smallskip
\renewcommand{\arraystretch}{1.12}%
\begin{tabular}{@{}r@{\enspace}l|l@{}}
\ixobs{Mai 13. Untere Kulm.}{319^\circ}{50'}{20''{,}73}{220}{5}{3{,}94}{\ixobsCount{3}{Beob.}}{\ixobsCount{4}{»}}\\[.6ex]
\ixobs{13. Obere Kulm.}{323}{8}{41{,}51}{216}{46}{44{,}31}{\ixobsCount{1}{»}}{\ixobsCount{1}{»}}
\end{tabular}
\end{center}

% Printed 41; principal PDF 48; Google PDF 52.
\par\clearpage
\noindent\makebox[12mm][l]{}\makebox[\dimexpr\linewidth-24mm\relax][c]{\small\textbf{III.\quad ERGEBNISSE.}}\makebox[12mm][r]{\fontsize{13}{16}\selectfont 41}\par\medskip

\begin{center}
1824.\quad Kreis im Osten.\par\smallskip
\renewcommand{\arraystretch}{1.18}%
\begin{tabular}{@{}r@{\enspace}l|l@{}}
\ixobs{Apr, 20. Obere Kulm.}{323^\circ}{7'}{52''{,}62}{216}{48}{54{,}93}{\ixobsCount{1}{Beob.}}{\ixobsCount{2}{»}}\\[.7ex]
\ixobs{21. Untere Kulm.}{319}{52}{30{,}27}{220}{4}{19{,}32}{\ixobsCount{3}{»}}{\ixobsCount{4}{»}}\\[.7ex]
\ixobs{21. Obere Kulm.}{323}{7}{54{,}16}{216}{48}{54{,}21}{\ixobsCount{3}{»}}{\ixobsCount{4}{»}}\\[.7ex]
\ixobs{25. Untere Kulm.}{319}{52}{30{,}03}{220}{4}{21{,}10}{\ixobsCount{3}{»}}{\ixobsCount{4}{»}}\\[.7ex]
\ixobs{27. Obere Kulm.}{323}{7}{55{,}70}{216}{48}{52{,}93}{\ixobsCount{3}{»}}{\ixobsCount{4}{»}}\\[.7ex]
\ixobs{28. Obere Kulm.}{323}{7}{55{,}40}{216}{48}{52{,}22}{\ixobsCount{3}{»}}{\ixobsCount{4}{»}}\\[.7ex]
\ixobs{29. Untere Kulm.}{319}{52}{29{,}17}{220}{4}{21{,}34}{\ixobsCount{3}{»}}{\ixobsCount{4}{»}}\\[.7ex]
\ixobs{Mai 1. Untere Kulm.}{319}{52}{28{,}59}{220}{4}{22{,}62}{\ixobsCount{3}{»}}{\ixobsCount{4}{»}}\\[.7ex]
\ixobs{1. Obere Kulm.}{323}{7}{57{,}22}{216}{48}{51{,}66}{\ixobsCount{3}{»}}{\ixobsCount{4}{»}}\\
\multicolumn{2}{@{}c|}{\rule{0pt}{3.3ex}Kreis im Westen.}&\\[.5ex]
\ixobs{Mai 2. Untere Kulm.}{40^\circ}{4'}{20''{,}00}{139}{52}{27{,}15}{\ixobsCount{3}{Beob.}}{\ixobsCount{4}{»}}\\[.7ex]
\ixobs{8. Obere Kulm.}{36}{48}{49{,}32}{143}{7}{57{,}63}{\ixobsCount{3}{»}}{\ixobsCount{4}{»}}\\[.7ex]
\ixobs{9. Untere Kulm.}{40}{4}{22{,}93}{139}{52}{25{,}68}{\ixobsCount{3}{»}}{\ixobsCount{4}{»}}
\end{tabular}
\end{center}

Die Änderungen der Deklination des Polarsterns ergeben sich aus \textsc{Bessel}s
Tafeln wie folgt:
\par\smallskip
\noindent\textbf{IX.}\hfill 6\hspace{.12\textwidth}\par
\par
\endgroup
